Three setup mistakes that silently cost accuracy, made concrete.

\examplenote{\code{examples/py/ex09\_pitfalls.py}}
\exfile{examples/py/ex09_pitfalls.py}{Python}
\resultfile{results/ex09_pitfalls_py.json}{Measured pitfall deltas + status.}

\begin{enumerate}
\item \textbf{Wrong ITRF realization} --- a measured ${\sim}7\,$mm LEO
divergence over one hour (ITRF2014 vs ITRF2020). Real but second-order;
match your reference data's realization.
\item \textbf{Missing EOP} --- check
\code{get\_init\_status().eop\_loaded}. Without it UT1$-$UTC defaults to 0
$\Rightarrow$ ${\sim}3\,$m LEO error. Run \code{Loader().load\_eop()}.
\item \textbf{Stale leap seconds} --- embedded coverage ends 2025;
querying later dates without a refreshed file gives seconds of TAI/UTC
error $\Rightarrow$ kilometers of position.
\end{enumerate}

\pitfall{Forces are presence-as-switch, identical in both languages: SRP
and drag default \emph{OFF}, and \code{srp=SrpConfig(...)} /
\code{drag\_env=...} turn them on (the old \code{enable\_srp} flag --- and
its C++/Python default divergence --- is gone). They default off only
because they need the satellite's mass, areas and (for drag) space weather
--- \emph{not} because a POE/SP3 reference ``already includes'' them: the
reference fixes only the start state, the forces after it must still be
modeled. Turning them on lowers 1-hour RMSE for every 2025 Sentinel (SRP
${\sim}{-}28\%$; drag $-17\%$ at 700\,km); \code{recommended\_config()}
does both in one call.}
